\documentclass[11pt]{article}
% ===========================================================================
%  Shared preamble for the AP Calculus AB exam-preparation worksheets.
%  Input right after \documentclass. Files beginning with "_" are not compiled
%  on their own. Navy and gold, to match the course's lesson pages.
% ===========================================================================
\usepackage[letterpaper,top=2.3cm,bottom=2.2cm,left=1.7cm,right=1.7cm,headheight=28pt,footskip=30pt]{geometry}
\usepackage{amsmath,amssymb}
\usepackage[T1]{fontenc}
\usepackage[utf8]{inputenc}
\usepackage{lmodern}
\usepackage{xcolor}
\usepackage[most]{tcolorbox}
\usepackage{enumitem}
\usepackage{fancyhdr}
\usepackage{array}
\usepackage{tikz}
\usepackage{pgfplots}
\pgfplotsset{compat=1.18}

\definecolor{navy}{HTML}{161E54}
\definecolor{ink}{HTML}{27357E}
\definecolor{gold}{HTML}{B98A00}
\definecolor{goldbg}{HTML}{FFFDF3}
\definecolor{inkbg}{HTML}{F6F7FD}
\definecolor{rule}{HTML}{C9CFF0}
\definecolor{okgreen}{HTML}{1B7A44}
\definecolor{okbg}{HTML}{F1F8F3}

\pagestyle{fancy}
\fancyhf{}
\fancyhead[L]{\footnotesize Name: \underline{\hspace{4cm}}}
\fancyhead[C]{\textbf{\color{navy}Integration Academy}\\[-2pt]{\scriptsize Calculus AB \textperiodcentered{} AP\textsuperscript{\textregistered} Exam Preparation}}
\fancyfoot[L]{\tiny\color{black!55}\parbox[t]{0.78\textwidth}{AP\textsuperscript{\textregistered} is a registered trademark of the College Board, which was not involved in the production of, and does not endorse, this material. \textcopyright{} Integration Academy.}}
\fancyfoot[R]{\scriptsize\thepage}
\renewcommand{\headrulewidth}{1.2pt}
\renewcommand{\headrule}{\hbox to\headwidth{\color{gold}\leaders\hrule height \headrulewidth\hfill}}
\renewcommand{\footrulewidth}{0pt}
\setlength{\parindent}{0pt}
\setlength{\parskip}{4pt}

% Title block: \wstitle{code}{title}{unit line}
\newcommand{\wstitle}[3]{\begin{center}{\small\bfseries\color{gold}WORKSHEET #1}\\[3pt]{\LARGE\bfseries\color{navy}#2}\\[3pt]{\small\color{black!60}#3}\end{center}\vspace{2pt}}
% Section heading
\newcommand{\sech}[1]{\par\vspace{10pt}{\large\bfseries\color{navy}#1}\par\nobreak\vspace{2pt}{\color{rule}\hrule height0.8pt}\vspace{7pt}}
% A level tag printed beside a group of questions
\newcommand{\level}[1]{\par\vspace{6pt}{\footnotesize\bfseries\color{gold}\MakeUppercase{#1}}\par\nobreak\vspace{2pt}}

\newtcolorbox{ideabox}{colframe=ink,colback=inkbg,coltitle=white,colbacktitle=ink,fonttitle=\bfseries\small,title={Key ideas},arc=2pt,boxrule=0.8pt,left=8pt,right=8pt}
\newtcolorbox{example}[1]{colframe=rule,colback=white,coltitle=navy,colbacktitle=inkbg,fonttitle=\bfseries\small,title={#1},arc=2pt,boxrule=0.8pt,borderline west={3pt}{0pt}{ink},left=10pt,right=8pt}
\newtcolorbox{frbox}[1]{colframe=gold,colback=goldbg,coltitle=white,colbacktitle=gold,fonttitle=\bfseries\small,title={#1},arc=2pt,boxrule=0.8pt,breakable,left=8pt,right=8pt}
\newtcolorbox{solbox}{colframe=okgreen,colback=okbg,coltitle=white,colbacktitle=okgreen,fonttitle=\bfseries\small,title={Worked solutions},arc=2pt,boxrule=0.8pt,breakable,left=8pt,right=8pt}
\newtcolorbox{scorebox}[1]{colframe=navy,colback=white,coltitle=white,colbacktitle=navy,fonttitle=\bfseries\small,title={#1},arc=2pt,boxrule=0.8pt,breakable,left=8pt,right=8pt}

\newcommand{\soln}{\par\smallskip\textbf{\color{ink}Solution.}\ }
\newcommand{\result}[1]{\par\smallskip\textbf{#1}}
% blank working space of a given height
\newcommand{\work}[1]{\par\vspace{#1}}
% a point value in the margin of a scoring guide
\newcommand{\pts}[1]{\hfill{\footnotesize\bfseries\color{gold}[#1]}}

% Coordinate grid for graph questions: \begin{apgraph}{xmin}{xmax}{ymin}{ymax} ... \end{apgraph}
\newenvironment{apgraph}[4]{\begin{center}\begin{tikzpicture}\begin{axis}[
  axis lines=middle, xmin=#1, xmax=#2, ymin=#3, ymax=#4,
  grid=both, grid style={black!12}, xtick distance=1, ytick distance=1,
  ticklabel style={font=\scriptsize}, xlabel={$x$}, ylabel={$y$},
  width=10cm, height=7.2cm, samples=80, clip=false,
  every axis plot/.append style={line width=1.1pt, color=ink}]}{\end{axis}\end{tikzpicture}\end{center}}
\newcommand{\closeddot}[2]{\addplot[only marks, mark=*, mark size=2.2pt, color=ink] coordinates {(#1,#2)};}
\newcommand{\opendot}[2]{\addplot[only marks, mark=*, mark size=2.2pt, color=ink, mark options={fill=white}] coordinates {(#1,#2)};}

% Multiple-choice options: four across, or stacked when the options are long.
\newcommand{\opts}[4]{\par\vspace{3pt}\begin{tabular}{@{}p{0.235\linewidth}p{0.235\linewidth}p{0.235\linewidth}p{0.235\linewidth}@{}}(A)\ #1 & (B)\ #2 & (C)\ #3 & (D)\ #4\end{tabular}\par\vspace{9pt}}
\newcommand{\optsv}[4]{\par\vspace{3pt}\begin{tabular}{@{}l@{\ \ }p{0.9\linewidth}@{}}(A) & #1\\(B) & #2\\(C) & #3\\(D) & #4\end{tabular}\par\vspace{9pt}}

\fancyhead[R]{\footnotesize\color{black!60}Worksheet 1.5}
\begin{document}

\wstitle{1.5}{Resolving 0/0: Factoring, Conjugates and Compound Fractions}{Unit 1: Limits and Continuity}

\begin{ideabox}
\begin{itemize}[leftmargin=*,itemsep=2pt,topsep=2pt]
  \item $\tfrac00$ is \textbf{indeterminate}: a signal to rewrite, never an answer.
  \item Polynomials: factor and cancel. Square roots: multiply by the conjugate. Fractions within fractions: use a common denominator.
  \item Absolute values and piecewise rules: take one-sided limits separately.
  \item Keep \textquotedblleft{}$\lim$\textquotedblright{} on every line that still contains the variable; drop it on the line where you substitute.
\end{itemize}
\end{ideabox}

\sech{Worked examples}
\begin{example}{Example 1 \textperiodcentered{} Factoring}
Evaluate $\displaystyle\lim_{x\to 2}\dfrac{x^2+3x-10}{x^2-4}$.\soln $\displaystyle\lim_{x\to 2}\dfrac{(x+5)(x-2)}{(x-2)(x+2)}=\displaystyle\lim_{x\to 2}\dfrac{x+5}{x+2}$.\result{The limit is $\dfrac74$.}
\end{example}

\begin{example}{Example 2 \textperiodcentered{} Conjugate}
Evaluate $\displaystyle\lim_{x\to 0}\dfrac{\sqrt{x+25}-5}{x}$.\soln Multiply by $\dfrac{\sqrt{x+25}+5}{\sqrt{x+25}+5}$: $\dfrac{x}{x\left(\sqrt{x+25}+5\right)}=\dfrac{1}{\sqrt{x+25}+5}$.\result{The limit is $\dfrac1{10}$.}
\end{example}

\begin{example}{Example 3 \textperiodcentered{} Compound fraction}
Evaluate $\displaystyle\lim_{x\to 3}\dfrac{\frac1x-\frac13}{x-3}$.\soln $\dfrac1x-\dfrac13=\dfrac{3-x}{3x}$, so the expression is $\dfrac{3-x}{3x(x-3)}=-\dfrac{1}{3x}$.\result{The limit is $-\dfrac19$.}
\end{example}

\newpage
\sech{Practice}

\level{Core}
\begin{enumerate}[leftmargin=*,itemsep=2pt]
  \item Evaluate $\displaystyle\lim_{x\to -1}\dfrac{x^2-1}{x+1}$.\work{2.3cm}
  \item Evaluate $\displaystyle\lim_{x\to 3}\dfrac{x^2-2x-3}{x-3}$.\work{2.3cm}
  \item Evaluate $\displaystyle\lim_{h\to 0}\dfrac{(4+h)^2-16}{h}$.\work{2.3cm}
  \item Evaluate $\displaystyle\lim_{x\to 1}\dfrac{x^3-1}{x-1}$.\work{2.3cm}
\end{enumerate}

\level{Exam level}
\begin{enumerate}[leftmargin=*,itemsep=2pt,start=5]
  \item Evaluate $\displaystyle\lim_{x\to 25}\dfrac{\sqrt x-5}{x-25}$.\work{2.8cm}
  \item Evaluate $\displaystyle\lim_{x\to 1}\dfrac{\frac{1}{x+2}-\frac13}{x-1}$.\work{2.8cm}
  \item Find the value of $k$ for which $\displaystyle\lim_{x\to -2}\dfrac{x^2+kx+6}{x+2}$ exists, and find the limit.\work{2.8cm}
  \item Evaluate $\displaystyle\lim_{x\to 5^-}\dfrac{|x-5|}{x^2-25}$.\work{2.8cm}
\end{enumerate}

\begin{frbox}{Question 9 \textperiodcentered{} Free-response practice}
Let $f(x)=\dfrac{x^2-9}{\sqrt{x+1}-2}$ for $x\neq3$, $x\ge-1$.
\begin{enumerate}[label=(\alph*),leftmargin=*,itemsep=2pt]
  \item Show that direct substitution at $x=3$ gives an indeterminate form.\work{2.6cm}
  \item Evaluate $\displaystyle\lim_{x\to 3}f(x)$, showing correct limit notation at each step.\work{2.6cm}
  \item What value must be assigned to $f(3)$ to make $f$ continuous at $x=3$? Give a reason.\work{2.6cm}
\end{enumerate}
\end{frbox}

\newpage
\begin{solbox}
\begin{enumerate}[leftmargin=*,itemsep=4pt]
  \item $x-1\to-2$.
  \item $\dfrac{(x-3)(x+1)}{x-3}=x+1\to4$.
  \item $\dfrac{8h+h^2}{h}=8+h\to8$.
  \item $x^2+x+1\to3$.
  \item $\dfrac{\sqrt x-5}{(\sqrt x-5)(\sqrt x+5)}=\dfrac{1}{\sqrt x+5}\to\dfrac1{10}$.
  \item $\dfrac{3-(x+2)}{3(x+2)(x-1)}=\dfrac{-(x-1)}{3(x+2)(x-1)}=-\dfrac{1}{3(x+2)}\to-\dfrac19$.
  \item The numerator must vanish at $-2$: $4-2k+6=0$, so $k=5$. Then $x^2+5x+6=(x+2)(x+3)$ and the limit is $1$.
  \item For $x<5$, $|x-5|=-(x-5)$, so the expression is $-\dfrac{1}{x+5}\to-\dfrac1{10}$.
  \item \begin{enumerate}[label=(\alph*),leftmargin=*,itemsep=3pt]
    \item Numerator: $9-9=0$. Denominator: $\sqrt4-2=0$. The form is $\tfrac00$.
    \item $\displaystyle\lim_{x\to 3}\dfrac{(x^2-9)\left(\sqrt{x+1}+2\right)}{(x+1)-4}=\displaystyle\lim_{x\to 3}\dfrac{(x-3)(x+3)\left(\sqrt{x+1}+2\right)}{x-3}=\displaystyle\lim_{x\to 3}(x+3)\left(\sqrt{x+1}+2\right)=6\cdot4=24$.
    \item $f(3)=24$, because continuity at 3 requires $f(3)=\displaystyle\lim_{x\to 3}f(x)$.
  \end{enumerate}
\end{enumerate}
\end{solbox}

\end{document}
